\section{Future Directions}
\label{sec:future}

To make the agenda actionable rather than aspirational, we anchor it on four measurable
quantities, defined with concrete testbeds in Table~\ref{tab:future_metrics}: the
\emph{success-vs-interactions curve} (does competence rise with autonomous experience?), the
\emph{skill-library reuse rate} (does stored competence compound?), the \emph{cross-embodiment
transfer drop} (how much is lost when a skill changes bodies?), and the \emph{provenance check}
(can a distributed skill's origin and test record be verified?). The first two instantiate
naturally on LIBERO-style lifelong task streams~\citep{libero}; the third on held-out-embodiment
splits of Open X-Embodiment~\citep{openx}; the fourth on marketplace
listings~\citep{unistore2026}.

% Future-agenda metrics: metric / definition / testbed / family stressed.
\begin{table*}[t]
\centering
\footnotesize
\setlength{\tabcolsep}{4pt}
\renewcommand{\arraystretch}{1.35}
\begin{tabular}{@{}p{0.17\textwidth} p{0.30\textwidth} p{0.27\textwidth} p{0.17\textwidth}@{}}
\toprule
\textbf{Metric} & \textbf{Definition} & \textbf{Instantiation (testbed)} & \textbf{Family stressed} \\
\midrule
Success-vs-interactions curve &
held-out success rate as a function of accumulated autonomous interactions; report the curve
and its area, not a single endpoint &
lifelong task streams in the style of LIBERO-100~\citep{libero}, fixed interaction budget &
\S\ref{sec:code} full loop; \S\ref{sec:vla} fine-tuning \\
Skill-library reuse rate &
fraction of new-task solutions that invoke at least one previously stored skill, and mean
invocations per stored skill &
growing code libraries~\citep{voyager} evaluated across LIBERO-style task suites~\citep{libero} &
\S\ref{sec:skill} libraries; \S\ref{sec:code} memory rung \\
Cross-embodiment transfer drop &
difference in success between the training embodiment and an unseen one, at matched task and
budget &
held-out-embodiment splits of Open X-Embodiment~\citep{openx} (22 embodiments) &
\S\ref{sec:vla}; \S\ref{sec:transfer} \\
Provenance check &
fraction of distributed skills carrying a verifiable origin, training-data statement, and
reproducible test record &
audits of marketplace listings~\citep{unistore2026} against a declared provenance schema &
market skills (\S\ref{sec:open}) \\
\bottomrule
\end{tabular}
\caption{\textbf{An actionable protocol for the future agenda}: four measurable quantities, their
definitions, and concrete testbeds. The first two instantiate LIBERO-style lifelong
learning~\citep{libero}; the third instantiates Open X-Embodiment-style multi-robot
transfer~\citep{openx}; the fourth targets the skill marketplaces of \S\ref{sec:open}.}
\label{tab:future_metrics}
\end{table*}


Three directions follow directly from the map, each now with its success measure.
\textbf{(1) Standardised evaluation of self-improvement.} The systems on the top rung
(\S3.1.5) each report improvement on their own splits; the field lacks a shared protocol that
plots the success-vs-interactions curve and the skill-library reuse rate under a common interaction
budget. A benchmark that publishes both curves per system, in the way LIBERO standardised
lifelong manipulation~\citep{libero}, would make ``self-improvement'' a reported number rather
than a claim. \textbf{(2) Adaptation as a first-class marketplace primitive.} The open problems
of \S\ref{sec:open} suggest a research programme in which a downloaded skill is not a frozen
artefact but a starting point that the \S\ref{sec:code} loop specialises to the local robot and
environment; its success measures are the transfer drop before versus after on-device
adaptation, and the provenance check on what was downloaded. \textbf{(3) Bridging the RL and
code flavours of skills.} \S\ref{sec:skill} shows that skills are learned either as latent RL
policies or as retrievable code; a unifying interface, code that wraps, calls, and refines
learned policies, and learned policies distilled back into inspectable code, would let the
self-improvement machinery of \S\ref{sec:code} operate over both, and its progress is legible in
the reuse rate crossing family boundaries. We see the combination of a standard self-improvement
benchmark with an adaptation-centric marketplace as the most likely path from today's static
skill stores to genuinely capable, compounding robots.
